\scanpage{84}
% archive foliation: 83
\margin{$\unsure{c}$}

$\mathcal M$ catégorie motivique \inserted{\unsure{$k$ un corps}},
de gerbe $\mathcal G$. Supposons que le lien $\mathcal L$ de $\mathcal G$
soit \uline{commutatif}, donc correspond à un groupe algébrique
\inserted{comm de type multiplicatif} $G$ sur $\mathbb Q$.
Ainsi, $G^\circ$ est un tore. Soit $E$ un motif ``fidèle''.
\deleted{et \unsure{par $E$ un foncteur}}
\inserted{Alors} $G$ opère sur $E$, considéré comme
\deleted{\unsure{objet}}\inserted{section}
de la catégorie fibrée $\operatorname{HOM}(\mathcal G,\uuline{\unsure{\mathrm{Proj}}})$,
et en particulier $G(\mathbb Q)$ opère sur $E$ donc
\[
G(\mathbb Q)\overset{\varphi}{\hookrightarrow}\operatorname{End}_{\mathcal M}(E).
\]
D'ailleurs,
\deleted{les éléments de $G(\mathbb Q)$}
\deleted{Soit $U=\operatorname{Ker}(G^\circ\longrightarrow G_m)$,
alors $U$ opère}
\deleted{\unsure{sur}}
Choisissons une forme de polarisation sur $E$, d'où une structure
d'Algèbre de Riemann sur $A=\operatorname{End}_{\mathcal M}(E)$,
et \unsure{on trouve} que pour $u\in G(\mathbb Q)$,
\[
\varphi(u)\varphi(u)^{\deleted{\unsure{*}}\inserted{'}}=
\varepsilon(u)\unsure{1|1_E}\deleted{\unsure{*}}.
\]
[En particulier, si $u\in\operatorname{Ker}\varepsilon$,
\inserted{$U=$}
on aura $\varphi(u)\varphi(u)'=\operatorname{id}$.]
Donc l'algèbre engendrée par les $\varphi(u)$ dans $A$ est auto-adjointe,
et $U(\mathbb Q)$ apparaît comme un s-groupe du
``groupe unitaire'' de $A$, formé des unités $u$ telles que $uu'=1$.
\deleted{Plus précisément}
Or ce \inserted{dernier} groupe est
\inserted{\unsure{fermé dans $\operatorname{Units}(A/\mathbb Q)$} et}
\unsure{un} \uline{donc} \uline{compact}.
On \unsure{trouve} que $U^\circ$ est un \uline{tore compact}.
